\section{Experimental Setup}
\label{sec:agents}

All methods are evaluated on the same accepted L1-L5 instances and scored by the same contract from Section~\ref{sec:scoring}. For each level, every compared policy is run on the same paired accepted-seed manifest, so later differences are measured on matched causal worlds rather than on different sampled instances. The comparisons in this section differ only in which part of the active-discovery pipeline is supplied, fixed, or left to the policy itself.

\subsection{Ablation design}

The compared policy set separates three questions over the same benchmark family: how well a method constructs observational structure from $D^{\mathrm{obs}}$; how well it uses interventions once a reasonable observational graph is available; and how well it integrates both stages end to end. The hybrid policies make those cuts explicit. Per-policy allowed-action sets, prompts, and code paths are pinned down in Appendix~\ref{app:policies:perpolicy}.

The classical reference policy, PC + greedy (\texttt{pc\_greedy}), applies PC with Fisher-$z$ conditional-independence tests at $\alpha=0.05$ to the observational sample \citep{spirtes2000causation}. It then resolves the remaining undirected edges with a fixed intervention heuristic that intervenes on the highest-undirected-degree node and orients adjacent edges with a $z$-calibrated mean-shift rule (Appendix~\ref{app:policies:zrule}).

The two end-to-end LLM policies keep the full active task exposed. LLM raw (\texttt{llm\_raw}) receives the observational matrix, may intervene, and must submit the final mixed graph itself. LLM stats (\texttt{llm\_stats}) keeps that same burden but adds Fisher-$z$-style \texttt{correlation}, \texttt{partial\_correlation}, and \texttt{independence\_test} actions (Appendix~\ref{app:policies:stats}), so the comparison isolates whether explicit CI-style statistical primitives improve active discovery under the same runtime.

The two hybrid policies remove one stage of the pipeline at a time. PC CPDAG + LLM (\texttt{pc\_cpdag\_llm}) supplies the PC partial graph together with observational summary statistics, but withholds the raw observational matrix; what remains is the active use of interventions and the final orientation of the graph. LLM stats + greedy (\texttt{llm\_stats\_cpdag\_greedy}) makes the complementary cut: the model uses the statistical interface to construct a partial graph and submit it, after which the same greedy intervention resolver completes the active stage.

\subsection{Model panel}

The paper evaluates a compact five-model panel: GPT-5.5, GPT-5.4-mini, Sonnet-4.6, Haiku-4.5, and Gemini-3-Flash. The panel is small by design: it is large enough to test whether the same decomposition remains visible across a meaningful capability spread, but not so large that the paper turns into a model ranking exercise.

\subsection{Run protocol}

The main paper reports only L1-L5. For each level, all policies and models are evaluated on the same accepted manifest of eight seeds (Appendix~\ref{app:repro:seeds}), so all reported comparisons are paired at the instance level.

All LLM policies run in the same sequential single-action loop (Appendix~\ref{app:policies:loop}): at each step the model receives the current session state and the accumulated interaction history, then emits exactly one action from its policy-specific action set. Parallel tool use is disabled, and the final allowed step is submit-only. Raw-style policies use a maximum of $20$ actions per episode, while stats-style policies use a maximum of $40$.
